\originalchapter{鞅、可选停止与风险中性定价}
\sourcelecture{34--36}
本章同时使用官方补充讲义.

条件期望描述给定信息后的平均判断. 鞅把这种判断沿时间排列: 今天的数值等于考虑今天全部信息后对明天数值的条件期望. 这一性质比“各时刻均值相同”强得多. 本章先研究公平赌博与停止规则, 再在可直接计算的市场模型中建立风险中性概率, 最后推导欧洲期权的 Black--Scholes 公式. 金融部分始终讨论给定假设下的数学定价, 不把市场价格当作客观事件频率.

\originalsection{信息、过滤与鞅}
在掷硬币实验中, 时刻 $n$ 已知前 $n$ 次结果, 但未知下一次结果. 为准确表达“此刻知道什么”, 把所有此刻可以判定真假的事件收集起来.

\begin{definition}[过滤与适应过程]
在概率空间 $(\Omega,\mathcal F,\Prob)$ 上, 一个 $\sigma$ 代数是含 $\Omega$、对补集和可数并集封闭的事件集合. 过滤是一列递增的 $\sigma$ 代数
\[
 \mathcal F_0\subseteq\mathcal F_1\subseteq\cdots\subseteq\mathcal F.
\]
$\mathcal F_n$ 表示时刻 $n$ 的信息. 若每个 $X_n$ 都是 $\mathcal F_n$ 可测的, 即对每个实数 $a$, $\{X_n\leq a\}\in\mathcal F_n$, 则过程 $(X_n)$ 称为适应该过滤. 由 $X_0,\ldots,X_n$ 产生的最小 $\sigma$ 代数记为 $\sigma(X_0,\ldots,X_n)$, 给出自然过滤.
\end{definition}
可测性在这里的含义是 $X_n$ 的值可以由当前信息确定. 若另有一组观察 $Y_0,\ldots,Y_n$, 过滤可以包含它们, 不能默认只有价格本身的历史. 同一个过程在不同信息集合下, 可能具有不同的鞅性质.

在有限信息树上, $\E(Z\mid\mathcal F_n)$ 的定义很具体: 对每个当前已知的历史节点, 用条件概率在其可能后续结果上取平均. 一般可积随机变量的条件期望是一个 $\mathcal F_n$ 可测可积变量 $W$, 满足
\[
 \E[W\ind_A]=\E[Z\ind_A]\quad(A\in\mathcal F_n).
\]
它在几乎处处意义下唯一. 有限树和离散随机游走的实例可按有限或可数历史逐项条件化; 连续取值的增量及布朗运动实例则使用第8章已声明的一般条件期望存在性工具, 不在这里重新构造该定理. 使用的性质是线性、已知有界因子可提出、独立变量的条件期望等于其均值, 以及塔式法则
\[
 \E[\E(Z\mid\mathcal F_{n+1})\mid\mathcal F_n]=\E(Z\mid\mathcal F_n).
\]
塔式法则可从上述定义核对: 左边对任何 $A\in\mathcal F_n$ 的积分等于 $\E[Z\ind_A]$, 且对 $\mathcal F_n$ 可测, 因而由唯一性得到等式. 在有限树上这正是先按后一天分组平均, 再按今天分组平均等于直接按今天平均.

\begin{definition}[鞅]
适应过程 $(M_n)_{n\geq0}$ 若满足每个 $n$ 上 $\E|M_n|<\infty$, 且
\[
 \E(M_{n+1}\mid\mathcal F_n)=M_n\quad\text{几乎必然},
\]
则称为相对于 $(\mathcal F_n)$ 和 $\Prob$ 的鞅. 若等号换为 $\geq$ 或 $\leq$, 则分别称为次鞅或上鞅.
\end{definition}
可积性保证条件期望有意义. 适应性保证右侧是当前已知的量. 条件均值条件则要求对每一种具有正概率的当前信息都公平, 不能只要求整体平均公平. 对鞅反复使用塔式法则得
\[
 \E(M_m\mid\mathcal F_n)=M_n\quad(m\geq n),\qquad
 \E M_n=\E M_0.
\]
第一个公式用 $m-n$ 归纳, 第二个公式再取无条件期望.

\begin{example}
设 $A_1,A_2,\ldots$ 独立可积, $\mathcal F_n=\sigma(A_1,\ldots,A_n)$. 证明下列鞅构造: 零均值增量之和, 单位均值因子之积, 单位方差增量的平方补偿, 以及逐次更新的条件期望.
\end{example}
\begin{solution}
若 $\E A_i=0$, 则 $S_n=\sum_{i=1}^nA_i$ 适应且可积, 并且
\[
 \E(S_{n+1}\mid\mathcal F_n)=S_n+\E A_{n+1}=S_n.
\]
若改为 $\E A_i=1$, 则 $P_n=\prod_{i=1}^nA_i$, $P_0=1$, 也是鞅. 独立性给出 $\E|P_n|=\prod_{i=1}^n\E|A_i|<\infty$, 而
$\E(P_{n+1}\mid\mathcal F_n)=P_n\E A_{n+1}=P_n$.
这包括每步独立乘以 $1+a$ 或 $1-a$、两者各概率 $1/2$ 的价格模型, $0<a<1$.

若 $\E A_i=0$, $\E A_i^2=1$, 则 $S_n^2-n$ 可积, 且
\begin{align*}
 \E[S_{n+1}^2-(n+1)\mid\mathcal F_n]
 &=S_n^2+2S_n\E A_{n+1}+\E A_{n+1}^2-(n+1)\\
 &=S_n^2-n.
\end{align*}
其中 $S_nA_{n+1}$ 可积由独立性和两者可积保证. 最后若 $Z$ 可积, $M_n=\E(Z\mid\mathcal F_n)$ 适应且
$\E|M_n|\leq\E|Z|$, 塔式法则给出
$\E(M_{n+1}\mid\mathcal F_n)=M_n$.
\end{solution}

\begin{example}
$A$ 以相同概率取 $1,-1$, 令 $X_0=0$, $X_n=(-1)^nA$ ($n\geq1$), 使用自然过滤. 判定它是否为鞅.
\end{example}
\begin{solution}
每个 $X_n$ 均值为零, 但观察 $X_1$ 后便确定 $A$. 所以 $n\geq1$ 时下一步已经完全可预测,
\[
 \E(X_{n+1}\mid\mathcal F_n)=-X_n\ne X_n.
\]
故不是鞅. 即使随机游走本来是鞅, 若把所有未来掷币结果在时刻零就告诉观察者, 它也不再对扩大后的过滤是鞅.
\end{solution}

若 $Z=\ind_A$, 条件期望鞅就是 $\Prob(A\mid\mathcal F_n)$. 例如观察一场比赛逐分钟更新获胜概率, 只要所有更新来自同一个事先选定的概率模型, 就形成鞅. 一条具体轨迹可以不断上升或下降; 鞅的等式是在每个历史下对尚未看到的信息取平均. 若事先概率模型估计不准, 该模型下仍可有鞅, 所以“概率更新是鞅”本身不证明估计符合现实.

\originalsection{停止时间与停止过程}
\begin{definition}[停止时间]
随机变量 $T:\Omega\to\{0,1,2,\ldots\}\cup\{\infty\}$ 称为停止时间, 若对每个 $n\geq0$,
\[
 \{T\leq n\}\in\mathcal F_n.
\]
即到时刻 $n$ 时能判断是否已经停止.
\end{definition}
这等价于 $\{T=n\}\in\mathcal F_n$ 对每个 $n$ 成立. 一方向使用
$\{T=n\}=\{T\leq n\}\setminus\{T\leq n-1\}$ 和过滤递增性; 在 $n=0$ 时前一个事件视为空集. 另一方向使用
$\{T\leq n\}=\bigcup_{k=0}^n\{T=k\}$.

首次到达集合 $B$ 的时间 $T=\inf\{n:X_n\in B\}$ 是停止时间, 因为
$\{T\leq n\}=\bigcup_{k=0}^n\{X_k\in B\}$ 属于 $\mathcal F_n$. 第 $j$ 次访问 $B$ 的时间也一样, 因为截至 $n$ 的访问次数 $\sum_{k=0}^n\ind_{\{X_k\in B\}}$ 可由当前历史计算. 这里要求 $B$ 为可测集合, 空集合的首次到达时间约定为 $\infty$.

“未来十步内的最后一次到达零”通常不是停止时间. 对简单对称随机游走, 定义 $L=\max\{0\leq n\leq2:S_n=0\}$. 同样的时刻零信息下, 两条后续路径 $++$ 与 $+-$ 分别使 $L=0$ 与 $L=2$, 所以 $\{L=0\}$ 不属于平凡的 $\mathcal F_0$. “截至固定期限取得最大值的时刻”也需要看未来, 一般不是停止时间. 在某个特定退化过程中它们可能是, 因而判断应看事件判据, 不能只靠规则名称.

\begin{proposition}[停止后的过程仍是鞅]
若 $M$ 是鞅, $T$ 是停止时间, 则 $M_{n\wedge T}$ 也是鞅. 其中 $n\wedge T=\min(n,T)$.
\end{proposition}
\begin{proof}
$M_{n\wedge T}=\sum_{k=0}^{n-1}M_k\ind_{\{T=k\}}+M_n\ind_{\{T\geq n\}}$, 所以适应且可积. 一步增量为
\[
 M_{(n+1)\wedge T}-M_{n\wedge T}
 =\ind_{\{T>n\}}(M_{n+1}-M_n).
\]
$\{T>n\}\in\mathcal F_n$, 因而取条件期望并将已知指示因子提出, 得到零. 这正是停止过程的鞅条件.
\end{proof}
停止之后把数值冻结, 于是每一步都只在“尚未停止”时参与下一次公平变化. 是否参与由变化发生之前的信息决定, 所以一步平均仍为零.

\originalsection{可选停止定理及其边界}
\begin{theorem}[有界停止时间]
若 $M$ 为可积鞅, $T$ 为停止时间, 且 $T\leq N$ 几乎必然, 其中 $N$ 是确定的有限整数, 则 $M_T$ 可积且
\[
 \E M_T=\E M_0.
\]
\end{theorem}
\begin{proof}
可积性由 $|M_T|\leq\sum_{k=0}^N|M_k|$ 保证. 路径上有恒等式
\[
 M_T=M_0+\sum_{k=0}^{N-1}\ind_{\{T>k\}}(M_{k+1}-M_k).
\]
对每项, 指示变量由时刻 $k$ 的信息决定, 所以
\[
 \E[\ind_{\{T>k\}}(M_{k+1}-M_k)]
 =\E\left[\ind_{\{T>k\}}\E(M_{k+1}-M_k\mid\mathcal F_k)\right]=0.
\]
有限求和后得到结论. 这也揭示了为何必须用停止时间: 若是否参与下一步取决于下一步输赢, 指示变量不能从条件期望提出.
\end{proof}

\begin{theorem}[有界过程与几乎必然有限的停止时间]
若 $M$ 为鞅, $T<\infty$ 几乎必然, 且存在确定常数 $C$ 使
$|M_{n\wedge T}|\leq C$ 对所有 $n$ 几乎必然成立, 则
$\E M_T=\E M_0$. 特别地, 一致有界鞅满足此条件.
\end{theorem}
\begin{proof}
$T\wedge N$ 是有界停止时间, 所以前一定理给出
$\E M_{T\wedge N}=\E M_0$. 在 $\{T\leq N\}$ 上,
$M_{T\wedge N}=M_T$; 在其余事件上两者的差绝对值至多 $2C$. 因而
\[
 |\E M_{T\wedge N}-\E M_T|\leq2C\Prob(T>N)\longrightarrow0.
\]
最后一个极限使用 $\{T>N\}$ 递减到 $\{T=\infty\}$, 后者概率为零. 这给出期待等式, 无需把极限交换当作无条件操作.
\end{proof}
这里明确要求 $T$ 几乎必然有限, 否则 $M_T$ 未定义, 除非先证明并定义无穷时刻的极限. 所谓鞅有界是所有时刻共用一个确定界, 不是“对每个固定 $n$, $M_n$ 都有界”. 简单随机游走在每个固定 $n$ 都满足 $|S_n|\leq n$, 仍不满足一致有界条件.

\begin{proposition}[可预测下注]
设 $H_k$ 是 $\mathcal F_k$ 可测有界变量. 若 $M$ 是鞅, 则截至固定期限 $N$ 的净收益
\[
 G_N=\sum_{k=0}^{N-1}H_k(M_{k+1}-M_k)
\]
具有期望零; $G_n$ 本身是鞅.
\end{proposition}
\begin{proof}
有限和可积. 每项取条件期望为
$H_k\E(M_{k+1}-M_k\mid\mathcal F_k)=0$. 一步增量同样满足条件均值为零, 因而收益过程是鞅.
\end{proof}
下注大小可以依赖已发生的输赢, 但不能偷看本次结果. 这一定理允许在固定期限内采用相当复杂的策略, 不把无穷期、无资金上限的策略自动包括进来.

\originalsection{赌场破产、往返访问与无界停止反例}
设整数 $i,b$ 满足 $0<i<b$, $S_0=i$, 每步独立以概率 $1/2$ 加 $1$ 或减 $1$. 定义
$T=\inf\{n:S_n\in\{0,b\}\}$. 这是首次到达边界的停止时间.

首先证明 $T<\infty$ 几乎必然. 任意尚未退出的状态, 接下来连续 $b$ 次向上都必然触及上边界, 该串的概率为 $2^{-b}$. 将时间分成长度 $b$ 的块, 每块给定前史后仍有至少 $2^{-b}$ 的退出概率, 因而
\[
 \Prob(T>kb)\leq(1-2^{-b})^k\longrightarrow0.
\]
这个界虽不精确, 但足以支持下面所有停止论证.

\begin{theorem}[公平赌场破产概率与平均时长]
上述随机游走满足
\[
 \Prob(S_T=b)=\frac ib,\qquad \Prob(S_T=0)=1-\frac ib,
 \qquad \E T=i(b-i).
\]
\end{theorem}
\begin{proof}
停止过程 $S_{n\wedge T}$ 值在 $[0,b]$ 内, 是有界鞅, 所以
$\E S_T=i$. 若成功概率为 $p$, 则 $bp=i$, 得 $p=i/b$.

由平方补偿例子, $S_n^2-n$ 是鞅. 对有界停止时间 $T\wedge N$ 使用可选停止,
\[
 \E[S_{T\wedge N}^2]-\E(T\wedge N)=i^2.
\]
左边第一项因 $0\leq S_{T\wedge N}\leq b$ 及 $T<\infty$ 而趋于
$\E S_T^2=b^2(i/b)=bi$. 具体误差至多 $b^2\Prob(T>N)$.
又 $\E(T\wedge N)=\sum_{k=0}^{N-1}\Prob(T>k)$, 非负级数的部分和趋于 $\E T$ (可由整数变量的尾和公式直接定义). 所以
$\E T=bi-i^2=i(b-i)$, 同时证明了平均时长有限.
\end{proof}
在一般区间 $[a,b]$ 内, 若从 $c$ 出发的鞅几乎必然在有限时刻无跳越地触及某一端点, 且停止过程一致有界, 则先触及 $b$ 的概率为 $(c-a)/(b-a)$. 若过程可能越过端点, 终值未必恰为 $a,b$, 不能照搬这条线性公式.

\begin{example}
公平随机游走从 $6$ 出发, 遇到 $0$ 或 $12$ 便结束. 求它在结束前依次达到 $4,8,4,8$ 的概率.
\end{example}
\begin{solution}
先从 $6$ 在 $[4,12]$ 内求先达 $4$ 的概率, 为 $(12-6)/(12-4)=3/4$. 达到 $4$ 后, 在 $[0,8]$ 内先达 $8$ 的概率为 $4/8=1/2$. 从 $8$ 再在 $[4,12]$ 内先达 $4$ 的概率为 $1/2$. 最后由 $4$ 再先达 $8$ 的概率为 $1/2$. 总概率为 $3/32$.
这里每次到达后未来掷币仍独立公平. 为避免把这个事实仅作为定理名引用, 可对每个可能到达时间 $n$ 条件化: 该到达事件只由前 $n$ 次掷币决定, 而任一固定长的后续掷币串与这前史独立. 后续串的概率不随 $n$ 或具体前史改变, 再对 $n$ 求和, 就得到每次重启时相同的退出公式和概率乘积.
\end{solution}

\begin{example}
从 $S_0=0$ 出发的简单对称随机游走, 令 $T=\inf\{n:S_n=1\}$. 证明 $T<\infty$ 几乎必然, 但 $\E S_T\ne\E S_0$.
\end{example}
\begin{solution}
对每个正整数 $a$, 在 $[-a,1]$ 内先到 $1$ 的概率为 $a/(a+1)$. 这一事件包含于 $\{T<\infty\}$, 所以
$\Prob(T<\infty)\geq a/(a+1)$ 对每个 $a$ 成立, 令 $a\to\infty$ 得概率为 $1$. 然而 $S_T=1$, 所以 $\E S_T=1\ne0$.
每个 $T\wedge N$ 仍满足 $\E S_{T\wedge N}=0$. 虽然
$S_{T\wedge N}\to1$ 几乎必然, 期望却不收敛到 $1$: 未停止的少数路径可能已经走到很负的值, 它们维持整体均值为零. 这里不存在共同的确定界控制所有停止值.
\end{solution}

另一反例是输后加倍. 第一次下注 $1$, 每次输后下注额翻倍, 首次赢后停止. 前 $n$ 次全输的概率为 $2^{-n}$, 损失为 $2^n-1$; 只要前 $n$ 次内赢过一次, 收益为 $1$. 因而截至 $n$ 的期望为
\[
 (1-2^{-n})\cdot1+2^{-n}\bigl(-(2^n-1)\bigr)=0.
\]
无限策略几乎必然最终赢且收益 $1$, 但最大资金需求没有界. 这与有界期限策略的零期望不矛盾, 也表明不能仅凭几乎必然有限的停止时间就交换期望和极限.

\originalsection{无套利与一期二叉树}
现在以股票和无风险账户建立一个完全可以求解的市场. 股票现价 $S>0$, 一期后只能为 $uS$ 或 $dS$, 其中 $0<d<u$. 无风险账户一期增长因子为 $R=1+r>0$. 可以买卖任意实数股, 允许借贷和卖空, 同一资产买卖价格相同, 无交易费用、税费或违约, 股票不分红. 两种状态都具有正的实际概率. 这些是假设, 不由概率公理推导.

\begin{definition}[套利与复制]
一期套利是初始净投入为零、两种终值均非负、至少一种终值严格为正的组合. 若一个组合在每一种状态的终值均与某合约相同, 称为该合约的复制组合.
\end{definition}
初始负成本且终值非负的策略也能变为零成本套利: 将收到的钱存入无风险账户. 反之若两个相同终值的组合价格不同, 买便宜的、卖贵的, 再存入差价, 即构成套利. 所以无套利迫使复制组合与合约价格相同.

\begin{theorem}[一期无套利条件]
上述市场无套利当且仅当
\[
 d<R<u.
\]
此时唯一的风险中性向上概率为
\[
 q=\frac{R-d}{u-d}\in(0,1),\qquad qu+(1-q)d=R.
\]
\end{theorem}
\begin{proof}
若 $R\leq d$, 借入 $S$ 买一股, 初始成本零, 终值为
$S(u-R)$ 或 $S(d-R)$, 均非负且上状态严格正, 是套利.
若 $R\geq u$, 卖空一股把所得 $S$ 存入账户, 终值为
$S(R-u)$ 或 $S(R-d)$, 同样是套利.

反之若 $d<R<u$, 上述 $q$ 在 $(0,1)$. 任一组合由 $\Delta$ 股股票及金额 $B$ 的账户组成, 初价 $V_0=\Delta S+B$, 两终值为
$V_u=\Delta uS+RB$, $V_d=\Delta dS+RB$. 由 $qu+(1-q)d=R$,
\[
 qV_u+(1-q)V_d=R(\Delta S+B)=RV_0.
\]
若 $V_0=0$, 且终值非负并至少一个严格正, 左侧由于权重均正便严格正, 与右侧零矛盾. 所以无套利. 解均值方程给出的 $q$ 唯一.
\end{proof}
这里实际向上概率 $p$ 没有出现在公式里. $q$ 的作用是把价格写成贴现均值, 不要求 $q=p$. 无套利禁止无风险的确定盈利, 允许带风险的正期望投资; 因而“实际概率下有正期望收益”不等同于套利.

\begin{theorem}[复制与一期定价]
合约在上涨和下跌时分别支付 $H_u,H_d$. 其唯一复制组合为
\[
 \Delta=\frac{H_u-H_d}{S(u-d)},\qquad
 B=\frac{uH_d-dH_u}{R(u-d)}.
\]
无套利价格为
\[
 V_0=\Delta S+B=\frac1R\bigl(qH_u+(1-q)H_d\bigr).
\]
\end{theorem}
\begin{proof}
要求 $\Delta uS+RB=H_u$, $\Delta dS+RB=H_d$. 两式相减得到 $\Delta$, 代回得 $B$. 因 $S(u-d)\ne0$, 线性方程唯一可解. 对两终值用 $q,1-q$ 加权, 得 $RV_0=qH_u+(1-q)H_d$. 若合约报价偏离 $V_0$, 同终值组合的买卖差价就给出套利, 故价格唯一.
\end{proof}

\begin{example}
股票现价 $80$, 一期后为 $100$ 或 $60$, 账户增长因子为 $1.05$. 执行价为 $85$ 的看涨合约终值为 $(S_1-85)^+$. 求复制组合和价格.
\end{example}
\begin{solution}
$u=1.25$, $d=0.75$, $q=(1.05-0.75)/0.5=0.6$. 两种支付为 $15,0$. 因而
\[
 \Delta=\frac{15}{40}=\frac38,\qquad
 B=\frac{-0.75(15)}{1.05(0.5)}=-\frac{150}7\approx-21.4286.
\]
购买 $3/8$ 股并借入 $150/7$, 初始净投入 $30-150/7=60/7$. 上涨终值 $37.5-22.5=15$, 下跌终值 $22.5-22.5=0$. 同样,
$V_0=(0.6\cdot15)/1.05=60/7$.
\end{solution}

\originalsection{多期复制与贴现鞅}
把一期模型重复 $N$ 次. 每个历史节点的下一步仍为当前价乘以 $u$ 或 $d$, 两个分支的实际条件概率均为正, 无风险增长因子仍为 $R$, 且 $d<R<u$. 在完整的 $2^N$ 条涨跌路径上, 定义新概率 $Q$: 每次上涨概率为 $q$, 下跌为 $1-q$, 独立. 实际概率与 $Q$ 因而具有相同的零概率事件. 对一条有 $k$ 次上涨的具体路径, $Q$ 概率为 $q^k(1-q)^{N-k}$; 同一终价对应多条路径时应加上路径数, 不能把它们当作单一路径.

\begin{proposition}
在 $Q$ 下, 贴现股价 $S_n/R^n$ 是鞅. 任意终期支付 $H$ (可依赖完整路径) 的价格为
\[
 V_n=R^{-(N-n)}\E_Q(H\mid\mathcal F_n),\qquad
 V_0=R^{-N}\E_QH.
\]
此价格可以由只使用当前信息的股票和账户交易逐期复制, 且 $V_n/R^n$ 为 $Q$ 鞅.
\end{proposition}
\begin{proof}
每个节点上
$\E_Q(S_{n+1}\mid\mathcal F_n)=S_n[qu+(1-q)d]=RS_n$, 所以贴现过程为鞅. 从终期令 $V_N=H$, 逐层向后定义
\[
 V_n=\frac1R(qV_{n+1}^{u}+(1-q)V_{n+1}^{d}).
\]
在该节点应用一期复制公式, 选择
\[
 \Delta_n=\frac{V_{n+1}^{u}-V_{n+1}^{d}}{S_n(u-d)},\qquad
 B_n=V_n-\Delta_n S_n.
\]
下一步组合恰有价值 $V_{n+1}$. 到达该下一节点时, 将组合重新调整成该节点所需的股票和账户持仓; 调整前后总价值同为 $V_{n+1}$, 不需要外部注资, 这就是自融资. 归纳直到 $N$, 便复制 $H$.

逐层条件化或塔式法则把递归展开为
$V_n=R^{-(N-n)}\E_Q(H\mid\mathcal F_n)$. 因而
$V_n/R^n=R^{-N}\E_Q(H\mid\mathcal F_n)$ 是条件期望鞅. 同终值复制组合的无套利论证逐期成立, 故得到价格唯一.
\end{proof}
这段论证实际构造了概率和策略, 没有调用一般金融基本定理. 完整二叉树中任何终期支付都能复制, 称为市场完备. 如果每个节点有三个可能股票价而只有股票和账户两个交易工具, 一般有三个终值方程但仅两个未知持仓, 不能保证所有合约都可复制, 风险中性概率也可能不唯一.

\begin{example}
股票 $S_0=100$, 两期内每期乘 $1.2$ 或 $0.8$, 账户 $R=1$. 求执行价 $100$ 的两期看涨价格及第一期持仓.
\end{example}
\begin{solution}
$q=1/2$. 终价为 $144,96,96,64$, 终值为 $44,0,0,0$. 因而 $V_0=44/4=11$. 一期上涨节点价格为 $22$, 下跌节点为 $0$. 初始持股
$\Delta_0=(22-0)/(120-80)=0.55$, 账户金额 $B_0=11-55=-44$.
若上涨到 $120$, 组合价值为 $22$, 下一步持股
$\Delta_1^u=44/(144-96)=11/12$, 账户金额 $22-(11/12)120=-88$.
若下跌到 $80$, 合约已无可能取得正终值, 该节点两种子节点支付均为零, 所以持股和账户金额均为零. 调整没有新增现金, 所有路径末端终值都正确.
\end{solution}

\originalsection{状态价格、风险中性概率与计价单位}
二叉树展示了风险中性概率如何由复制产生. 还有一个一般的有限状态说明, 可以把市场里的二元事件合约与概率公理联系起来. 固定终期 $T$, 有限个状态为 $\omega_1,\ldots,\omega_m$, 假设每个支付 $\ind_{\{\omega_i\}}$ 的合约都可买卖且被履约, 其当前价格为 $\pi_i$. 无风险支付 $1$ 的现价为 $D=\ee^{-rT}$.

同支付同价格迫使 $\sum_i\pi_i=D$, 因为各状态合约合起来总支付 $1$. 无套利迫使 $\pi_i\geq0$, 否则买该合约收到钱, 存入账户后即得非负且正的终值. 若每个状态实际有正概率且零价合约可交易, 还要求 $\pi_i>0$, 否则免费买入就形成套利. 定义
\[
 Q(\{\omega_i\})=\frac{\pi_i}{D},\qquad
 Q(A)=\sum_{\omega_i\in A}\frac{\pi_i}{D}.
\]
它确为有限状态空间上的概率. 任意支付 $H$ 为各状态合约的线性组合, 所以
\[
 \operatorname{Price}(H)=\sum_i\pi_iH(\omega_i)
 =D\E_QH.
\]
对于无限状态, 单靠无套利的有限买卖只能推出有限可加性; 要从无限事件合约推出可数可加性还需价格对单调极限的连续性等条件. 本章有限树证明没有这项缺口; 后面的连续模型直接假设有一个可数可加的风险中性概率, 不把其一般存在性误说成已证.

风险中性概率与实际频率不同, 因为不同状态的一元货币对人的价值可能不同. 经济不景气时多收到一元可能更有用, 保险支付可能因此更贵. 这只是解释价格与频率可能分离的经济直觉, 不是从无套利单独推出的严格大小关系. 买卖价差、交易限制及履约风险若存在, 同终值组合的差价也未必能按上述理想方式套利.

计价单位也影响概率. 在上述有限状态模型中, 设一个正资产终值为 $A_i>0$, 当前价 $A_0=\sum_i\pi_iA_i$. 以该资产为计价单位, 支付“一单位资产且仅在状态 $i$ 发生时”的现价相对于 $A_0$ 为
\[
 Q^A(\{\omega_i\})=\frac{\pi_iA_i}{A_0}
 =\frac{Q(\{\omega_i\})A_i}{\E_QA_T}.
\]
这些权重相加为 $1$, 且
\[
 \operatorname{Price}(H)=A_0\E_{Q^A}\left(\frac{H}{A_T}\right).
\]
所以更换计价单位确实改变概率权重. 这个有限求和推导给出精确公式, 无需调用抽象测度变换定理. 在某状态终值较低的计价资产, 该状态在新概率中的相对权重也较低.

\originalsection{看涨函数、看跌函数与分布恢复}
先暂时不贴现金融价格. 对非负且可积随机变量 $X$, 定义
\[
 \mathcal C(K)=\E(X-K)^+,\qquad
 \mathcal P(K)=\E(K-X)^+,\qquad K\geq0,
\]
其中 $a^+=\max(a,0)$, $F(x)=\Prob(X\leq x)$. 这两个函数只由概率分布决定.

\begin{proposition}[尾积分与看涨函数]
\[
 \mathcal C(K)=\int_K^\infty(1-F(x))\dd x,\quad
 \mathcal P(K)=\int_0^K F(x)\dd x,\quad
 \mathcal C(K)-\mathcal P(K)=\E X-K.
\]
$\mathcal C$ 非负、递减且凸, $\mathcal C(0)=\E X$, $\mathcal C(K)\to0$ 当 $K\to\infty$. 它的右、左导数分别为
\[
 \mathcal C'_+(K)=-\Prob(X>K),\qquad
 \mathcal C'_-(K)=-\Prob(X\geq K)\quad(K>0).
\]
若 $F$ 连续, 则 $\mathcal C'(K)=F(K)-1$. 若 $X$ 具有在 $K$ 连续的密度 $f$, 则 $\mathcal C''(K)=f(K)$.
\end{proposition}
\begin{proof}
对每个确定 $x\geq0$,
$(x-K)^+=\int_K^\infty\ind_{\{x>t\}}\dd t$,
$(K-x)^+=\int_0^K\ind_{\{x<t\}}\dd t$.
非负函数可以交换积分次序; 可由矩形上的非负简单函数逐次逼近证明. 所以第一式成立. 第二式的被积函数是 $\Prob(X<t)$, 它与 $F(t)$ 仅在原子点不同; 正概率原子至多可数, 不改变积分, 得所写公式. 也可用
$(x-K)^+-(K-x)^+=x-K$ 直接得到平价恒等式.

每个固定 $x$ 上 $(x-K)^+$ 都递减且凸, 取期望保留这两性质. 又
\[
 0\leq\mathcal C(K)\leq\E[X\ind_{\{X>K\}}]\longrightarrow0,
\]
尾部期望趋零由 $X$ 可积保证. 从右侧差商
\[
 \frac{(X-K-h)^+-(X-K)^+}{h},\qquad h>0,
\]
绝对值不超过 $1$, 点态极限为 $-\ind_{\{X>K\}}$. 对这些绝对值不超过 $1$ 的差商可交换期望和极限, 得右导数. 左侧相同论证给出 $-\ind_{\{X\geq K\}}$. 在无原子点两者相同, 得到一阶导数. $F'=f$ 在密度连续处成立, 再求导得到二阶公式.
\end{proof}
原子处两侧斜率的差为 $\Prob(X=K)$, 所以折角记录离散概率质量. 不能对任意分布写一个普通函数 $f=\mathcal C''$; 有原子时应使用斜率跳跃或分布意义的导数.

若 $X=S_T$ 是风险中性终价, 欧洲看涨和看跌期权的当前价格分别为
\[
 C_0(K)=\ee^{-rT}\mathcal C_Q(K),\qquad
 P_0(K)=\ee^{-rT}\mathcal P_Q(K).
\]
因为无分红股票满足 $\ee^{-rT}\E_QS_T=S_0$, 得到看涨看跌平价
\[
 C_0(K)-P_0(K)=S_0-K\ee^{-rT}.
\]
连续分布下
\[
 Q(S_T>K)=-\ee^{rT}C_0'(K),\qquad
 f_Q(K)=\ee^{rT}C_0''(K).
\]
第一阶导数来自很窄执行价区间的看涨价差, 第二阶导数来自相邻执行价的二阶差分. 真实报价只在有限个执行价上存在, 所以只能估计导数, 还受到噪声、买卖价差与插值影响. 这些式子恢复风险中性分布, 不直接恢复实际频率.

\begin{example}
$X$ 在 $1,3$ 两点上概率分别为 $2/3,1/3$. 写出 $\mathcal C$, 并说明如何由斜率恢复质量.
\end{example}
\begin{solution}
\[
 \mathcal C(K)=
 \begin{cases}
 5/3-K,&0\leq K\leq1,\\
 1-K/3,&1\leq K\leq3,\\
 0,&K\geq3.
 \end{cases}
\]
斜率从 $-1$ 在 $K=1$ 跳到 $-1/3$, 跳幅 $2/3$; 在 $K=3$ 再跳到 $0$, 跳幅 $1/3$. 这恰是两点概率. 普通二阶导数在其余点为零, 所以忽略折角会错失整个分布.
\end{solution}

\originalsection{二叉树的对数正态极限}
Black--Scholes 模型的核心假设是风险中性终价对数正态. 为解释其来源, 将总时长 $T$ 划为 $n$ 段, $h=T/n$, 取
\[
 u_h=\ee^{\sigma\sqrt h},\qquad d_h=\ee^{-\sigma\sqrt h},
 \qquad R_h=\ee^{rh},\qquad \sigma>0.
\]
$h$ 充分小时 $d_h<R_h<u_h$, 因为对数上的严格不等式为
$-\sigma\sqrt h<rh<\sigma\sqrt h$. 风险中性向上概率为
\[
 q_h=\frac{\ee^{rh}-\ee^{-\sigma\sqrt h}}
 {\ee^{\sigma\sqrt h}-\ee^{-\sigma\sqrt h}}.
\]
设第 $j$ 步的对数增量为 $Y_{h,j}$, 以概率 $q_h$ 取 $\sigma\sqrt h$, 否则取 $-\sigma\sqrt h$, 各步独立. 则
$\ln S_T^{(n)}=\ln S_0+\sum_{j=1}^nY_{h,j}$.

\begin{proposition}[固定终期的正态极限]
上述模型中
\[
 \ln S_T^{(n)}\ \Longrightarrow\
 \Normal\left(\ln S_0+(r-\sigma^2/2)T,\,\sigma^2T\right).
\]
因此 $S_T^{(n)}$ 分布收敛到相应对数正态变量.
\end{proposition}
\begin{proof}
令 $a=\sigma\sqrt h$. 泰勒展开的余项在 $h\downarrow0$ 时统一控制, 得
\begin{align*}
 \ee^{rh}-\ee^{-a}&=a+(r-\sigma^2/2)h+O(h^{3/2}),\\
 \ee^a-\ee^{-a}&=2a+O(h^{3/2}),
\end{align*}
所以
\[
 q_h=\frac12+\frac{r-\sigma^2/2}{2\sigma}\sqrt h+O(h),
 \quad m_h:=\E_QY_{h,1}=(r-\sigma^2/2)h+O(h^{3/2}),
\]
而 $v_h:=\Var_QY_{h,1}=\sigma^2h-m_h^2=\sigma^2h+O(h^2)$.
于是 $nm_h\to(r-\sigma^2/2)T$, $nv_h\to\sigma^2T$.
这里每行的分布随 $n$ 改变, 不能直接当作固定分布的独立同分布中心极限定理. 我们给出特征函数计算来核实所需版本.

设 $Z_h=Y_{h,1}-m_h$. 对固定实数 $t$, 使用
$\ee^{itZ_h}=1+itZ_h-t^2Z_h^2/2+O(|t|^3|Z_h|^3)$.
因为 $\E Z_h=0$ 且 $|Z_h|\leq2\sigma\sqrt h$, 有
\[
 \E\ee^{itZ_h}=1-\frac{t^2v_h}2+O(h^{3/2}).
\]
独立性使中心化和的特征函数为此式的 $n$ 次幂. 利用
$\ln(1+w)=w+O(w^2)$ (复数 $w$ 近零), 其对数为
\[
 n\ln\E\ee^{itZ_h}
 =-\frac{nt^2v_h}2+O(nh^{3/2})+O(nh^2)
 \longrightarrow-\frac{t^2\sigma^2T}2.
\]
加上均值因子便趋于所陈述正态分布的特征函数. 这里使用特征函数连续性定理: 若特征函数点态趋于某个在零连续的概率分布的特征函数, 则相应分布弱收敛. 此定理是特征函数极限判据, 本章不重新证明一般傅里叶反演理论; 上面的行列变化、误差累积与极限参数均已直接计算. 指数函数连续, 所以指数变换的分布也收敛; 等价地在每个 $x>0$ 上对 $\ln x$ 比较分布函数即可.
\end{proof}
注意 $-\sigma^2/2$ 修正不能省略. 股价的风险中性均值增长率为 $r$, 但对数的均值增长率为 $r-\sigma^2/2$, 原因是对数凹性与指数的正态均值修正.

仅有分布收敛不足以保证看涨期权的期望收敛, 因为支付 $(S-K)^+$ 无界. 在这个特定树上还可建立足够的尾部控制. 一步二阶矩因子为
\[
 q_h\ee^{2a}+(1-q_h)\ee^{-2a}
 =1+(2r+\sigma^2)h+O(h^{3/2}),
\]
于是独立性给出
\[
 \E_Q[(S_T^{(n)})^2]
 =S_0^2\bigl[1+(2r+\sigma^2)h+O(h^{3/2})\bigr]^n,
\]
对所有充分大的 $n$ 一致有界. 对非负变量 $S$, 尾部估计
$\E[S\ind_{\{S>M\}}]\leq\E S^2/M$ 说明这些股价的尾部期望一致趋零. 将支付截断为
$\min\{(S-K)^+,M\}$. 用非负尾积分公式,
\[
 \E\min\{(S-K)^+,M\}=\int_0^M\Prob(S>K+t)\dd t.
\]
极限对数正态分布的分布函数处处连续, 所以每个尾概率都收敛; 被积函数又在有限区间 $[0,M]$ 上被 $1$ 控制, 可交换极限与积分. 这直接证明截断支付的期望收敛, 无需另引一般的弱收敛等价刻画. 截断误差至多 $\E[S\ind_{\{S>K+M\}}]$, 由上述二阶估计一致控制; 极限对数正态变量也有有限二阶矩. 先令 $n\to\infty$, 再令 $M\to\infty$, 得到看涨期望也收敛. 每个树上的贴现因子 $R_h^{-n}=\ee^{-rT}$ 相同, 所以树价格趋于下面的公式.

这是一种固定终期分布及价格的极限证明, 不是整条随机路径收敛或连续时间复制策略收敛的证明. 若要严格构造连续时间的布朗运动及随机积分, 需要更进一步的概率理论. 对本章的欧洲终期支付, 固定终期极限已经够用.

\originalsection{Black--Scholes 模型与公式}\label{sec:bs}
连续时间的理想模型假定无分红股票、恒定无风险利率 $r$、恒定波动率 $\sigma>0$, 无交易摩擦且可连续交易. 在风险中性概率 $Q$ 下, 布朗运动 $W_t$ 有连续路径, $W_0=0$, 不重叠时间段的增量独立, $W_t-W_s\sim\Normal(0,t-s)$. 这里取自然信息 $\mathcal F_s=\sigma(W_u:0\le u\le s)$; 不额外预先透露未来增量. 定义
\[
 S_t=S_0\exp\left((r-\sigma^2/2)t+\sigma W_t\right).
\]
仅为了终期定价, 可以不先构造完整过程, 而直接假定
\[
 N:=\ln S_T\sim\Normal(\mu,\sigma^2T),\qquad
 S_0=\ee^{-rT}\E_QS_T.
\]
后一个等式是模型的贴现价格条件. 正态指数矩 $\E\ee^N=\ee^{\mu+\sigma^2T/2}$ 给出
\[
 \mu=\ln S_0+(r-\sigma^2/2)T.
\]
因此无需另行猜测真实世界下股票上涨趋势来填写 $\mu$. 它在风险中性模型里已由 $S_0,r,\sigma,T$ 决定.

完整过程也确实满足贴现鞅条件: 令 $A_t=\ee^{rt}$, 对 $s<t$,
\[
 \frac{S_t}{A_t}=\frac{S_s}{A_s}
 \exp\left(-\frac{\sigma^2}2(t-s)+\sigma(W_t-W_s)\right).
\]
未来增量与当前信息独立, 正态指数矩使后一个因子的均值为 $1$, 因而
$\E_Q(S_t/A_t\mid\mathcal F_s)=S_s/A_s$. 本章是在 $Q$ 下直接定义过程, 并未从任意实际概率下的股票模型证明风险中性测度存在, 也未调用 Girsanov 定理.

\begin{theorem}[欧洲看涨期权公式]
设 $S_0>0$, $K>0$, $T>0$, $\sigma>0$. 在上述模型中, 到期支付 $(S_T-K)^+$ 的当前价为
\[
 C_0=S_0\Phi(d_1)-K\ee^{-rT}\Phi(d_2),
\]
其中 $\Phi(z)=\int_{-\infty}^z(2\pi)^{-1/2}\ee^{-x^2/2}\dd x$,
\[
 d_1=\frac{\ln(S_0/K)+(r+\sigma^2/2)T}{\sigma\sqrt T},
 \qquad d_2=d_1-\sigma\sqrt T.
\]
\end{theorem}
\begin{proof}
设 $v=\sigma^2T$, $\mu=\ln S_0+(r-\sigma^2/2)T$, $k=\ln K$. 将支付拆成两个截断期望,
\[
 C_0=\ee^{-rT}\left(\E_Q[\ee^N\ind_{\{N>k\}}]
                 -KQ(N>k)\right).
\]
先算概率项. 标准化给出
\[
 Q(N>k)=1-\Phi\left(\frac{k-\mu}{\sqrt v}\right)
 =\Phi\left(\frac{\mu-k}{\sqrt v}\right)=\Phi(d_2),
\]
最后使用标准正态关于零的对称性.

指数项的困难是密度之外多出 $\ee^x$. 在指数中配方,
\[
 x-\frac{(x-\mu)^2}{2v}
 =\mu+\frac v2-\frac{(x-\mu-v)^2}{2v}.
\]
因而
\begin{align*}
 \E_Q[\ee^N\ind_{\{N>k\}}]
 &=\frac1{\sqrt{2\pi v}}\int_k^\infty
   \exp\left(x-\frac{(x-\mu)^2}{2v}\right)\dd x\\
 &=\ee^{\mu+v/2}\frac1{\sqrt{2\pi v}}\int_k^\infty
   \exp\left(-\frac{(x-\mu-v)^2}{2v}\right)\dd x\\
 &=\ee^{\mu+v/2}\Phi\left(\frac{\mu+v-k}{\sqrt v}\right)
 =S_0\ee^{rT}\Phi(d_1).
\end{align*}
将两项代回得到公式. 这也直接验证支付期望有限, 因为 $0\leq(S_T-K)^+\leq S_T$ 且 $\E_QS_T=S_0\ee^{rT}$.
\end{proof}
$d_2$ 给出的 $\Phi(d_2)$ 是风险中性行权概率 $Q(S_T>K)$.
$\Phi(d_1)$ 通常不是这个概率: 它来自额外乘以股价后的截断期望. 若按前述计价单位公式改用股票作为正计价资产, 则它确可解释为新概率下的行权概率. 这两个正态参数相差 $\sigma\sqrt T$, 正是配方造成的均值平移.

由看涨看跌平价, 欧洲看跌价格为
\[
 P_0=K\ee^{-rT}\Phi(-d_2)-S_0\Phi(-d_1).
\]
边界参数不必强行代入除以 $\sigma\sqrt T$ 的表达式: $T=0$ 时价格是 $(S_0-K)^+$; $K=0$ 时看涨支付就是股票, 价格为 $S_0$; $\sigma=0$ 时终价确定为 $S_0\ee^{rT}$, 看涨价格为 $(S_0-K\ee^{-rT})^+$.

\begin{example}
设 $S_0=K=100$, $r=0$, $T=1$, $\sigma=0.2$. 求看涨和看跌价格, 以及风险中性行权概率.
\end{example}
\begin{solution}
$d_1=0.1$, $d_2=-0.1$. 因而
\[
 C_0=100[\Phi(0.1)-\Phi(-0.1)]
 =100[2\Phi(0.1)-1]\approx7.9656.
\]
平价给出 $P_0=C_0$, 因为现价等于执行价且无利息. 风险中性行权概率为
$\Phi(-0.1)\approx0.46017$, 小于 $1/2$. 尽管 $\E_QS_T=100$, 对数均值为 $\ln100-0.02$, 中位数低于 $100$, 少量较大上涨补足均值.
\end{solution}

\originalsection{隐含波动率与模型限制}
波动率 $\sigma$ 表示单位时间对数增量的标准差, $\sigma^2$ 是方差率, 两者不能混用. 给定 $S_0,K,r,T$ 和看涨报价, 把使公式匹配报价的 $\sigma$ 称为隐含波动率. 先看公式对 $\sigma$ 是否严格递增, 才能讨论反解的唯一性.

\begin{proposition}[波动率敏感度]
令 $\phi(z)=(2\pi)^{-1/2}\ee^{-z^2/2}$. 对 $T,K,\sigma>0$,
\[
 \frac{\partial C_0}{\partial\sigma}
 =S_0\phi(d_1)\sqrt T>0.
\]
因此报价位于 $((S_0-K\ee^{-rT})^+,S_0)$ 内时, 存在唯一的正隐含波动率.
\end{proposition}
\begin{proof}
由 $d_1-d_2=\sigma\sqrt T$ 及
$d_1+d_2=2[\ln(S_0/K)+rT]/(\sigma\sqrt T)$,
\[
 \frac{\phi(d_1)}{\phi(d_2)}
 =\exp\left(-\frac{d_1^2-d_2^2}{2}\right)
 =\exp[-\ln(S_0/K)-rT],
\]
所以 $S_0\phi(d_1)=K\ee^{-rT}\phi(d_2)$.
对价格公式求导, 两项相消后剩下
\[
 \frac{\partial C_0}{\partial\sigma}
 =S_0\phi(d_1)\left(\frac{\partial d_1}{\partial\sigma}
             -\frac{\partial d_2}{\partial\sigma}\right)
 =S_0\phi(d_1)\sqrt T.
\]
当 $\sigma\downarrow0$, 公式极限为 $(S_0-K\ee^{-rT})^+$, 包括 $S_0=K\ee^{-rT}$ 时极限零. 当 $\sigma\to\infty$, $d_1\to\infty$, $d_2\to-\infty$, 所以价格趋于 $S_0$. 连续性和严格单调性给出开区间内反解存在且唯一.
\end{proof}
在同一股票、同一期限的理想恒波动率模型中, 所有执行价应给出同一个隐含波动率. 若不同执行价反推出不同数值, 它们描述的是报价相对于模型的差异, 而不是证明存在多个实际恒波动率. 波动率微笑或偏斜可以反映更厚的尾部、跳跃或波动率随状态变化等现象.

二叉树极限依赖小而独立、尺度为 $\sqrt h$ 的对数增量及恒定参数. 日内或跨日涨跌幅可能变化, 大跳跃不能按小增量泰勒余项控制, 真实利率和未来波动率也可能随机. 这些变化会改变终期分布. 本章没有证明任意真实股票都对数正态, 没有证明实际概率等于风险中性概率, 也没有用历史标准差自动替代未来模型参数. 连续时间精确复制涉及随机积分与自融资策略的可行性; 本章提供的是每棵有限二叉树的精确复制和其欧洲价格极限, 并完整计算给定对数正态风险中性模型下的公式.

\originalsection{练习与解答}
\begin{exercise}
公平随机游走从 $3$ 出发, 首次到达 $0$ 或 $10$ 时停止. 求到达 $10$ 的概率、平均停止时间, 并说明为什么不能直接对无界随机游走调用有界过程版本.
\end{exercise}
\begin{solution}
成功概率为 $3/10$, 平均停止时间为 $3(10-3)=21$. 使用的有界鞅是停止后的过程 $S_{n\wedge T}$, 它的值在 $[0,10]$ 内; 原过程 $S_n$ 本身没有统一界. 平均时长用 $S_n^2-n$ 在有界停止时间 $T\wedge N$ 上计算, 随后分别控制平方项和非负时间项的极限, 不能说 $S_n^2-n$ 是有界鞅.
\end{solution}

\begin{exercise}
股票现价 $50$, 一期后为 $65$ 或 $40$, 账户一期增长因子 $1.02$. 合约上涨支付 $8$, 下跌支付 $3$. 求风险中性概率、复制持仓和当前价格.
\end{exercise}
\begin{solution}
$u=1.3$, $d=0.8$, $q=(1.02-0.8)/(1.3-0.8)=0.44$. 持股
$\Delta=(8-3)/(65-40)=1/5$. 账户终值由下状态为
$3-(1/5)40=-5$, 所以当前账户金额 $B=-5/1.02=-250/51$.
价格 $V_0=10-250/51=260/51\approx5.09804$.
风险中性计算也为
$[0.44(8)+0.56(3)]/1.02=5.2/1.02=260/51$.
实际上涨概率可以为 $0.7$, 但不会因此改变这组复制终值或无套利价格.
\end{solution}

\begin{exercise}
非负终价 $X$ 的风险中性分布连续, 无风险贴现因子为 $D$. 证明对 $0<K_1<K_2$, 看涨价差满足
\[
 D\Prob_Q(X>K_2)\leq
 \frac{C_0(K_1)-C_0(K_2)}{K_2-K_1}
 \leq D\Prob_Q(X>K_1).
\]
说明怎样由很窄的价差估计一个尾概率.
\end{exercise}
\begin{solution}
由尾积分公式, 中间量等于
$D(K_2-K_1)^{-1}\int_{K_1}^{K_2}\Prob_Q(X>x)\dd x$.
尾概率随 $x$ 递减, 所以区间平均介于两端之间. 当两端同时趋于连续点 $K$, 两端尾概率趋于同一个值, 夹逼给出 $D\Prob_Q(X>K)$. 因此窄价差除以执行价差, 再除以 $D$, 可估计尾概率. 若有原子, 左右逼近分别给出包含和不包含该点的尾概率.
\end{solution}

\begin{exercise}
在 Black--Scholes 模型中求只在 $S_T>K$ 时支付 $1$ 的二元合约价格, 并与看涨价格的执行价导数比较.
\end{exercise}
\begin{solution}
终价对数标准化给出 $Q(S_T>K)=\Phi(d_2)$, 因而二元合约价格为 $\ee^{-rT}\Phi(d_2)$.
从看涨函数的一阶导数公式知
$-\partial C_0/\partial K=\ee^{-rT}\Phi(d_2)$.
也可以直接对 Black--Scholes 公式求导: 两个密度项因
$S_0\phi(d_1)=K\ee^{-rT}\phi(d_2)$ 及
$\partial d_1/\partial K=\partial d_2/\partial K=-1/(K\sigma\sqrt T)$ 相消, 留下 $-\ee^{-rT}\Phi(d_2)$. 这将概率计算、复制定价的含义和从价格恢复分布联系起来.
\end{solution}
